% Capitulo: Os quatro estágios da IA no trabalho com dados
\chapter{Os quatro estágios da IA no trabalho com dados}
\label{cap:estagios-ia}

A adoção da IA no trabalho com dados não aconteceu de uma vez. Ela veio em estágios, cada
um dando à IA um pouco mais de autonomia e mudando um pouco mais o papel de quem trabalha
com ela. Primeiro, a IA era uma conversa numa janela à parte. Depois, passou a trabalhar
dentro das ferramentas de código. Em seguida, passou a construir, publicar e manter
aplicações inteiras. E agora começa a se tornar proativa: em vez de esperar que alguém
pergunte, ela acompanha os dados, decide dentro de limites combinados e avisa quando uma
decisão humana é necessária.

Este capítulo descreve esses quatro estágios (Figura~\ref{fig:estagios-ia}). Eles não são
uma escada que toda empresa sobe na mesma ordem, e convivem no mesmo time: a mesma pessoa
pode conversar com a IA de manhã e supervisionar um agente à tarde. Mas separá-los ajuda a
entender o que muda em cada um e, principalmente, o que cada estágio exige dos dados e das
pessoas. A ideia de níveis de automação não é nova. Os carros autônomos, por exemplo, são
classificados em níveis que dizem quem faz o quê, o sistema ou o motorista, em cada situação
\parencite{sae2021}. A pergunta que organiza este capítulo é a mesma.

\begin{figure}[htbp]
  \centering
  % Os quatro estagios da IA no trabalho com dados (cap. 6).
\begin{tikzpicture}[diagrama, x=1mm, y=1mm,
    degrau/.style={caixa, anchor=south west, minimum width=27mm, text width=24mm,
      minimum height=17mm, align=center, font=\scriptsize},
    humano/.style={rotulo, anchor=north, text width=27mm}]
  \node[degrau] (e1) at (0,0)  {\textbf{1. Conversa}\\[0.5mm]a IA responde numa janela à parte};
  \node[degrau] (e2) at (30,7) {\textbf{2. Código}\\[0.5mm]a IA trabalha dentro do editor e do repositório};
  \node[degrau] (e3) at (60,14) {\textbf{3. Aplicações}\\[0.5mm]a IA constrói, publica e mantém sistemas inteiros};
  \node[degrau, fill=fundo, line width=0.7pt] (e4) at (90,21)
    {\textbf{4. Proativa}\\[0.5mm]a IA monitora, decide nos limites e avisa};
  \node[humano] at (e1.south) [yshift=-1mm] {a pessoa pergunta, copia e integra};
  \node[humano] at (e2.south) [yshift=-1mm] {a pessoa dirige e revisa};
  \node[humano] at (e3.south) [yshift=-1mm] {a pessoa especifica e supervisiona};
  \node[humano] at (e4.south) [yshift=-1mm] {a pessoa define as regras e decide nas exceções};
  \draw[seta, draw=cinza] (0,-12) -- (117,-12);
  \node[rotulo, anchor=north] at (58.5,-12.5) {mais autonomia para a IA, mais importância para as regras e para a revisão};
\end{tikzpicture}

  \caption{Os quatro estágios da IA no trabalho com dados. A cada estágio, a IA ganha
    autonomia e o papel humano muda de fazer para definir regras e decidir nas exceções.}
  \label{fig:estagios-ia}
  \fonte{elaborado pelo autor.}
\end{figure}

\section{Estágio 1: a IA como conversa}
\label{sec:estagio-conversa}

O primeiro contato da maioria das pessoas com a IA generativa foi uma janela de conversa. A
pessoa escreve uma pergunta, recebe uma resposta e decide o que fazer com ela. Para quem
trabalha com dados, isso já foi uma mudança grande: pedir uma consulta, colar um erro e pedir
explicação, pedir ajuda para estruturar uma análise ou para escrever o e-mail com as
conclusões (capítulo~\ref{cap:ia}).

A característica que define esse estágio é o isolamento. A IA não vê o banco de dados, não
conhece o código do projeto nem os padrões do time. Tudo o que ela sabe é o que a pessoa
cola na conversa, e tudo o que ela produz precisa ser copiado, adaptado e integrado à mão.
A pessoa é a ponte entre a IA e o trabalho real.

Esse isolamento tem um lado bom: o controle é total, porque nada acontece sem que a pessoa
execute. E tem um lado ruim: o contexto é pobre, a integração é manual e cansativa, e é fácil
colar na conversa o que não deveria sair da empresa (\veja{sec:seguranca}). O estágio da
conversa continua útil, e não vai desaparecer, mas é o que menos aproveita o que a IA pode
fazer.

\section{Estágio 2: a IA conectada ao código}
\label{sec:estagio-codigo}

O segundo estágio levou a IA para dentro das ferramentas em que o trabalho acontece. Primeiro
vieram os assistentes integrados ao editor de código, que sugerem a próxima linha ou o
próximo bloco enquanto a pessoa escreve. Depois vieram os agentes de programação: assistentes
que leem o repositório inteiro, entendem a estrutura do projeto, executam comandos, rodam os
testes, corrigem o que falhou e preparam a alteração para revisão. A pessoa deixa de pedir
trechos de código e passa a pedir tarefas.

A diferença em relação ao estágio anterior é o contexto e a ação. O agente vê o código real,
os testes, as convenções do time, e consegue agir sobre eles. Uma tarefa como \enquote{crie
o modelo de receita por cliente seguindo o padrão das outras tabelas, com os testes de
unicidade e de frete} pode ser executada do começo ao fim, com a pessoa revisando o resultado.

Esse estágio trouxe também um vocabulário novo e um debate útil. Em 2025, a expressão
\emph{vibe coding} se popularizou para descrever um jeito de programar em que a pessoa
descreve o que quer, aceita o código gerado sem lê-lo e segue em frente enquanto parece
funcionar. Ela foi rapidamente distinguida do uso responsável dos assistentes: se a pessoa
revisou, testou e entendeu o código que a IA escreveu, isso não é \emph{vibe coding}, é
programar com um assistente \parencite{willison2025}. Para protótipos descartáveis, aceitar
sem ler pode ser aceitável. Para um modelo de dados que vai alimentar a decisão de alguém,
não é. A forma que tem ganhado espaço para o trabalho profissional costuma ser chamada de
engenharia com agentes: a IA executa, a pessoa dirige e verifica. É o que os
capítulos~\ref{cap:programar} e~\ref{cap:testes} defendem com outras palavras.

\section{Estágio 3: construir, publicar e gerenciar aplicações}
\label{sec:estagio-aplicacoes}

No terceiro estágio, a IA deixa de trabalhar só sobre partes de um projeto e passa a
construir aplicações inteiras: a interface, a lógica, o banco de dados, a publicação e, cada
vez mais, a manutenção. Uma pessoa descreve uma ferramenta interna, como um painel de
acompanhamento com alertas ou uma aplicação de conciliação de cadastros
(capítulo~\ref{cap:aplicacao-propria}), e em pouco tempo tem uma versão funcionando,
publicada e acessível ao time.

Para times de dados, isso muda a fronteira do que é possível. Ferramentas internas que antes
nunca entravam na fila de prioridades da TI, porque eram pequenas demais, agora podem ser
construídas pelo próprio time. Uma aplicação para coletar dados que hoje chegam por planilha,
uma tela para o steward revisar sugestões, um painel com uma lógica que a ferramenta de BI
não suporta: tudo isso ficou ao alcance.

O risco também mudou de lugar. Construir ficou fácil; manter continua difícil. Uma aplicação
publicada precisa de alguém que cuide de segurança, de permissões, de atualização de
dependências, de dúvidas dos usuários e de mudanças no negócio. Uma empresa que deixa cada
pessoa publicar a sua aplicação com ajuda da IA, sem nenhum critério, repete o problema dos
mil dashboards (\veja{sec:mil-dashboards}), agora com sistemas inteiros. As mesmas práticas
que o livro propõe para os dados valem aqui: versionamento, testes, um dono para cada
aplicação, e um critério claro do que pode ir para produção (Parte~\ref{parte:devops}).

\section{Estágio 4: a IA proativa}
\label{sec:estagio-proativa}

Nos três primeiros estágios, a IA espera ser chamada. Alguém faz uma pergunta, pede uma
tarefa ou descreve uma aplicação. O quarto estágio inverte essa relação. A IA passa a
acompanhar os dados continuamente, com as regras de negócio no contexto, percebe quando um
número mudou de forma relevante, analisa a situação e, conforme o caso, decide sozinha dentro
de limites combinados ou avisa quem precisa decidir, já com a análise pronta.

É uma mudança profunda na natureza do dado dentro da empresa. Durante toda a história contada
nos capítulos anteriores, o dado foi passivo: ficava no armazém, ou no painel, esperando que
alguém o consultasse, percebesse algo e decidisse. A crítica de Ackoff aos relatórios que
ninguém lia (\veja{sec:sad}) e o problema dos painéis abertos uma vez por semana
(capítulo~\ref{cap:frequencia}) são versões do mesmo limite: a informação existe, mas só gera
ação se uma pessoa estiver olhando no momento certo. Com a IA proativa, o dado deixa de
esperar.

Uma referência clássica da engenharia de automação ajuda a pensar esse estágio com cuidado.
\textcite{parasuraman2000} propõem que a automação pode atuar sobre quatro funções: obter a
informação, analisá-la, escolher a decisão e executar a ação, cada uma num nível diferente,
do totalmente manual ao totalmente automático. Os painéis tradicionais automatizam a primeira
função e parte da segunda. A IA proativa avança sobre a terceira e, em alguns casos, sobre a
quarta. A pergunta de projeto deixa de ser \enquote{a IA consegue?} e passa a ser \enquote{em
que nível cada função deve ficar, para cada tipo de decisão?}.

Uma proposta recente para agentes de IA organiza a resposta pelo papel que a pessoa assume:
operadora, colaboradora, consultora, aprovadora ou observadora. No nível do aprovador, o agente
age sozinho e só envolve a pessoa em situações de risco ou previamente combinadas; no do
observador, age com autonomia total e a pessoa apenas acompanha \parencite{feng2025}. Para
decisões de negócio, o nível do aprovador costuma ser o ponto de equilíbrio, e a forma de
chegar lá é parecida com o que as empresas já fazem com pessoas: definir alçadas. A
Tabela~\ref{tab:decisoes-proativa} mostra um exemplo de como esse desenho pode ficar.

\begin{table}[htbp]
  \centering
  \caption{Um desenho de alçadas para uma IA proativa.}
  \label{tab:decisoes-proativa}
  \small
  \begin{tabularx}{\textwidth}{@{}>{\raggedright\arraybackslash}p{0.22\textwidth}
      >{\raggedright\arraybackslash}X>{\raggedright\arraybackslash}p{0.2\textwidth}
      >{\raggedright\arraybackslash}p{0.2\textwidth}@{}}
    \toprule
    Tipo de situação & Exemplo & O que a IA faz & Papel humano \\
    \midrule
    Rotineira, baixo impacto & Repor um item de estoque dentro do limite combinado & Decide, executa e registra & Audita por amostragem \\
    \addlinespace
    Fora do padrão, impacto médio & Custo de frete muito acima do contratado numa rota & Analisa e recomenda & Aprova ou rejeita \\
    \addlinespace
    Alto impacto ou ambígua & Suspender as compras de um fornecedor & Avisa com a análise pronta & Decide \\
    \addlinespace
    Fora das regras conhecidas & Uma situação que nenhuma regra cobre & Avisa que não sabe o que fazer & Define a regra nova \\
    \bottomrule
  \end{tabularx}
  \fonte{elaborado pelo autor com base em \textcite{parasuraman2000} e \textcite{feng2025}.}
\end{table}

A última linha da tabela é talvez a mais importante. Uma IA proativa bem desenhada sabe
reconhecer quando a situação está fora das regras que recebeu, e nesse caso a ação correta é
avisar, e não improvisar. É a velha ideia de gestão por exceção: as pessoas deixam de
acompanhar tudo e passam a cuidar do que foge do normal.

\rascunho{
  Relato ou exemplo da minha experiência: um caso em que um alerta proativo (com ou sem IA)
  mudou uma decisão, ou um caso em que eu gostaria de ter tido esse tipo de alerta. Pode ser
  também uma visão de como isso se aplicaria a um dos domínios do livro (agricultura,
  finanças ou logística).
}

\section{O que cada estágio exige}
\label{sec:estagios-exigem}

Cada estágio exige mais dos dados e das pessoas do que o anterior. No estágio da conversa,
basta saber perguntar e conferir. No estágio do código, é preciso saber revisar, o que exige
saber programar (capítulo~\ref{cap:programar}). No estágio das aplicações, é preciso
disciplina de engenharia para que o que foi construído possa ser mantido.

O estágio proativo é o mais exigente, porque nele a IA age sobre os dados sem que alguém
esteja olhando a cada passo. Para isso funcionar, os dados precisam estar prontos de uma forma
que poucas empresas já estão: regras de negócio explícitas e escritas, métricas definidas uma
vez (\veja{sec:camada-semantica}), cadastros conciliados (capítulo~\ref{cap:master-data}),
testes que garantam que os números estão certos (\veja{sec:tdd}), atualização na frequência
certa e monitorada (capítulo~\ref{cap:frequencia}), e documentação que sirva de contexto
(capítulo~\ref{cap:contexto-ia}). Não é coincidência que essa lista resuma boa parte do resto
do livro. A IA proativa é o estágio em que todo o trabalho de base deixa de ser boa prática e
passa a ser condição.

Há também um risco humano, bem conhecido na pesquisa sobre automação. Quando um sistema
automático acerta quase sempre, as pessoas tendem a confiar demais nele e a prestar menos
atenção justamente nos casos em que ele erra \parencite{parasuraman2010}. Uma IA proativa que
pede aprovação dezenas de vezes por dia, com recomendações quase sempre certas, corre o risco
de ser aprovada no automático. O desenho das alçadas precisa levar isso em conta: poucas
aprovações, bem escolhidas, com a informação necessária para decidir, valem mais do que muitas
aprovações que viram rotina.

\begin{destaque}
No último estágio, o dado deixa de esperar ser analisado. A IA acompanha, decide dentro de
limites combinados e chama uma pessoa quando a decisão é dela. O trabalho humano passa a ser
definir as regras e cuidar das exceções.
\end{destaque}

Com este capítulo, termina a primeira parte do livro: a história de como chegamos até aqui e
de para onde a IA está levando o trabalho com dados. A próxima parte muda de pergunta: com
tudo isso acontecendo, o que vale a pena estudar, e em que ordem?

\begin{leituras}
  \leitura{parasuraman2000}{O modelo clássico de tipos e níveis de automação. Escrito para
    sistemas como aviação e controle industrial, serve muito bem para pensar a IA nas
    decisões de negócio.}
  \leitura{feng2025}{Uma proposta atual de níveis de autonomia para agentes de IA, organizada
    pelo papel que a pessoa assume em cada nível.}
  \leitura{parasuraman2010}{Para entender por que as pessoas passam a confiar demais em
    sistemas automáticos, e como desenhar a interação para evitar isso.}
\end{leituras}
